\section{评估、可靠性与安全：从 demo 到生产系统}

本期访谈没有把 Agent 的问题停留在“能不能做一个炫酷 demo”。苏煜多次把问题拉回到 reliability、speed、cost effectiveness 和 deployment。一个 agent demo 可以通过精心设计的环境、样例和人工提示展示“它会做”；但生产系统要回答的是：它在陌生任务上失败率多高，失败是否可恢复，成本是否可控，权限边界是否清楚，数据泄漏或错误操作如何被限制。

\begin{importantbox}{从 demo 到生产的四个门槛}
第一是任务成功率，不只是单步动作正确；第二是可恢复性，失败后能否定位错误并继续；第三是成本，长程任务的 token、工具调用和人工兜底是否经济；第四是安全，agent 能操作真实环境时，权限、审计和最坏情况都必须被设计进去。
\end{importantbox}

\subsection{为什么 benchmark 不够}

早期 web agent 或 computer-use benchmark 让研究者能比较系统，但它们往往不能完整覆盖真实部署。真实企业环境里有登录态、权限、历史数据、组织流程、遗留系统、审计要求和用户不完整指令。一个 agent 在标准 benchmark 上表现不错，不代表它能在银行、法务、医疗、研发工具链或企业后台里安全运行。

\noindent\begin{tabularx}{\textwidth}{p{0.20\textwidth}p{0.34\textwidth}X}
\toprule
评估维度 & 需要测什么 & 为什么重要 \\
\midrule
Task success & 长程任务最终是否完成 & 避免只优化单步点击或单次回答。 \\
Recovery & 出错后能否自诊断、回滚、重试 & 决定系统是否能无人值守运行。 \\
Cost & token、工具调用、等待时间、人工介入 & 决定 agent 是否能规模化部署。 \\
Safety & 权限边界、危险动作、隐私与审计 & 从 demo 进入真实业务的必要条件。 \\
Adaptation & 是否能从新环境和新反馈中学习 & 连接 continual learning 与 expert agent。 \\
\bottomrule
\end{tabularx}

\subsection{Security 与 Safety 的差别}

访谈中有一个容易被忽略的区分：许多 safety 问题归根结底是能力问题，因为 agent 像新手一样不懂环境，容易犯低级错误；但 security 更偏 worst-case scenario，需要专门方法。也就是说，提升能力可以减少很多粗心错误，却不能替代权限控制、沙盒、审计和红队。

\begin{warningbox}{能力提升不是安全设计}
一个更聪明的 agent 可能更少误点按钮，但也可能更有能力绕过限制。生产系统必须把能力提升与安全机制分开设计：能力负责完成任务，安全负责限制任务空间、记录行为、阻止不可接受的动作。
\end{warningbox}


\subsection{实践清单：做一套 Agent 系统前先问什么}

如果把本期内容转成工程 checklist，至少要问八个问题。第一，Agent 的边界是什么，谁能给它目标，谁能中止它。第二，它面对的环境是什么，环境状态如何被观测。第三，它能采取哪些动作，哪些动作必须禁止或需要人工确认。第四，任务成功如何被自动判定，哪些任务必须人工验收。第五，它失败后如何恢复，是否能回滚或生成审计记录。第六，它如何积累经验，经验是否会污染后续任务。第七，它的成本是否可预测。第八，它会在哪个产品入口被用户自然调用。

\begin{importantbox}{工程化判断}
一个 Agent 系统能不能上线，不取决于它能否在演示中完成一次惊艳任务，而取决于上述八个问题能否被稳定回答。这个清单也是后续阅读其他 Agent 访谈时的复用框架。
\end{importantbox}

\subsection{本章小结}

Agent 的生产化不只是模型问题。评估、恢复、成本、安全和环境适配共同决定它能否从 demo 走向真实系统。这也是为什么 continual learning 与 world model 会成为主线：它们最终要服务可靠性，而不是只服务更漂亮的演示。

